\section{Simultaneous fixed spaces}\label{sec:fixed}

The rank hypothesis enters the proof through the following finite-dimensional
statement.  A unitary matrix may occur in several blocks, with either exponent
$1$ or exponent $-1$.  Only the signed numbers of its occurrences matter for
the lower bound on the total fixed-space dimension.
The proof extends the single-matrix incidence argument of
\cite[Lemma~3.1]{OpenAIKervaire2026} by using rational cohomology and
nested nonzero minors to accommodate several unitary matrix variables.

\begin{theorem}[Simultaneous fixed spaces]\label{thm:fixed-space}
Let $r\geq 0$ and $n\geq 1$, and let $k_1,\ldots,k_r$ be positive integers.
For each $j\in\{1,\ldots,n\}$, let $\mathcal B_j$ be a finite ordered list
of pairs $(i,s)$ with $1\leq i\leq r$ and $s\in\{1,-1\}$; repetitions
are allowed.  Set
\[
 d_j=\sum_{(i,s)\in\mathcal B_j}k_i,
 \qquad
 b_{ij}=\sum_{\substack{(i',s)\in\mathcal B_j\\i'=i}}s,
 \qquad B=(b_{ij}).
\]
Suppose that $B$ has row rank $r$ over $\mathbb Q$.  For arbitrary
$P_j\in\mathrm U(d_j)$, define
\[
 D=\prod_{i=1}^r\mathrm U(k_i),
 \qquad
 F_j(X)=\operatorname{diag}_{(i,s)\in\mathcal B_j}(X_i^s)P_j
 \quad (X\in D).
\]
Then there is $X\in D$ such that
\begin{equation}\label{eq:fixed-space-bound}
 \sum_{j=1}^n\dim_{\mathbb C}\ker(F_j(X)-I_{d_j})
 \ \geq\ \sum_{i=1}^r k_i.
\end{equation}
Here $\mathrm U(0)$ and an empty product of groups are points, and the
fixed-space dimension in dimension zero is zero.
\end{theorem}

\begin{proof}
\emph{Choosing target dimensions.}
The case $r=0$ is immediate.  Assume $r>0$, and reorder the rows so that
$k_1\geq\cdots\geq k_r$.  We first allocate the desired fixed-space
dimensions among the maps $F_j$.  There are distinct columns
$j_1,\ldots,j_r$ such that
\begin{equation}\label{eq:fixed-nested-minors}
 \det B[1,\ldots,t\mid j_1,\ldots,j_t]\ne 0
 \qquad (1\leq t\leq r).
\end{equation}
Indeed, the first $t$ rows have rank $t$.  The columns chosen at the previous
step remain independent on these rows, since their restrictions to the
first $t-1$ rows are independent; they can therefore be extended by one
column.  Define
\[
 l_{j_t}=k_t\quad(1\leq t\leq r),
 \qquad l_j=0\quad\text{for unchosen columns}.
\]
These choices satisfy $l_j\leq d_j$.  To see this for $j=j_t$, the
nonzero minor in \eqref{eq:fixed-nested-minors} implies that $b_{i,j_t}\ne0$
for some $i\leq t$.  The list $\mathcal B_{j_t}$ thus contains a block
of size $k_i\geq k_t$.  Moreover,
\begin{equation}\label{eq:fixed-target-dimensions}
 \sum_j l_j=\sum_i k_i,
 \qquad \sum_j l_j^2=\sum_i k_i^2=\dim_{\mathbb R}D.
\end{equation}
We will obtain $\dim\ker(F_j(X)-I_{d_j})\geq l_j$ simultaneously.
The equality of dimensions in \eqref{eq:fixed-target-dimensions} makes
this an intersection problem of complementary dimensions.  We develop
the cohomology classes that detect the required intersection.

\smallskip\noindent
\emph{Primitive cohomology classes.}
All cohomology in this proof has coefficients in $\mathbb Q$.  The
classical computation of unitary-group cohomology
\cite[Proposition~9.1 and Sections~19--20]{Borel1953} gives generators
\begin{equation}\label{eq:fixed-unitary-cohomology}
 H^*(\mathrm U(d);\mathbb Q)
   =\Lambda(e_{1,d},\ldots,e_{d,d}),
 \qquad |e_{a,d}|=2a-1,
\end{equation}
where $\Lambda$ denotes the exterior algebra.  Compatibility means that
inclusion as a coordinate block, with the identity on its complement,
satisfies
\begin{equation}\label{eq:fixed-block-restriction}
 \iota_{k,d}^*e_{a,d}=
 \begin{cases}
 e_{a,k},&a\leq k,\\
 0,&a>k.
 \end{cases}
\end{equation}
We will choose the generators to satisfy this compatibility and to be
\emph{primitive}: if $\mu$ is group multiplication, then
\begin{equation}\label{eq:fixed-primitivity}
 \mu^*e_{a,d}=e_{a,d}\otimes1+1\otimes e_{a,d}.
\end{equation}
Thus these classes add under multiplication, just as degree-one
classes do on a torus.

For completeness, these properties follow together from the last-column
bundle
\[
 \mathrm U(d-1)\longrightarrow\mathrm U(d)
       \longrightarrow S^{2d-1}.
\]
Start with the degree-one generator of $\mathrm U(1)=S^1$.  For $d\geq2$,
the base is simply connected, so the coefficient system in the
multiplicative Serre spectral sequence is constant
\cite[Theorem~5.2 and Section~5.3]{McCleary2001}.  Its second page is
\[
 E_2^{p,q}=H^p(S^{2d-1};\mathbb Q)
              \otimes H^q(\mathrm U(d-1);\mathbb Q).
\]
Its only nonzero columns are $p=0,2d-1$.  The only possible differential
between them is $d_{2d-1}$, which sends fiber degree $q$ to degree
$q-2d+2$.  Each fiber generator has degree at most $2d-3$, so its
differential has negative target degree and vanishes.
The differential vanishes on all their products by the
product rule for the differential.  Consequently the spectral sequence
collapses.  Its edge maps show that restriction to $\mathrm U(d-1)$ is
an isomorphism in degrees below $2d-1$, and that the pullback of the sphere's
top class is nonzero.  Lift each earlier generator in its degree, and take
this sphere pullback as $e_{d,d}$.  Their exterior products give the basis
in \eqref{eq:fixed-unitary-cohomology}: they give that basis in the
associated graded algebra supplied by the spectral sequence, and their
squares vanish because they have odd degree and the coefficients are
rational.  This also proves compatibility with the standard inclusions.

To check primitivity, subtract the two summands on the right of
\eqref{eq:fixed-primitivity} from its left side.  The difference restricts
to zero on either factor at the identity, so every term in its K\"unneth
decomposition has positive degree in both factors.  Each such degree is
at most $2d-2$.  Restriction to
$\mathrm U(d-1)\times\mathrm U(d-1)$ is injective on all these
bidegrees, by the edge-map isomorphism just proved.  On this product the
difference vanishes: for $a<d$ this is the induction hypothesis, and for
$a=d$ it follows because $e_{d,d}$ is pulled back by the last-column
projection, which is constant on the fiber.
This proves \eqref{eq:fixed-primitivity}; the case $d=1$ is immediate.
Coordinate-block inclusions in other positions are conjugate to the
standard one, and conjugation is homotopic to the identity in the
connected group $\mathrm U(d)$.  Finally, inversion sends $e_{a,d}$
to $-e_{a,d}$, by pulling \eqref{eq:fixed-primitivity} back along
$X\mapsto(X,X^{-1})$.

\smallskip\noindent
\emph{The fixed-space incidence class.}
We next construct a class that detects a fixed space of dimension at
least $l$, for $0\leq l\leq d$.  Let
\[
 R(d,l)=\{(W,E): W\in\mathrm U(d),\ E\in
        \operatorname{Gr}_l(\mathbb C^d),\ W|_E=I_E\},
 \qquad p(W,E)=W,
\]
where $\operatorname{Gr}_l(\mathbb C^d)$ is the complex Grassmannian of
$l$-dimensional subspaces.  Over a subspace $E$, the possible $W$ are
precisely the unitary transformations of $E^\perp$, extended by the
identity on $E$.  Orthonormal frame charts therefore make $R(d,l)$ a
smooth bundle over the Grassmannian with fiber $\mathrm U(d-l)$.
It is compact and without boundary, and
\[
 \dim_{\mathbb R}R(d,l)=2l(d-l)+(d-l)^2=d^2-l^2.
\]
The complex orientation of the Grassmannian and an orientation of
$\mathrm U(d-l)$ orient this bundle: changes of frame act on the fiber
by conjugation and preserve orientation.  Choose orientations and let
\[
 \alpha(d,l)=\operatorname{PD}\bigl(p_*[R(d,l)]\bigr)
        \in H^{l^2}(\mathrm U(d);\mathbb Q).
\]
Here $\operatorname{PD}$ is Poincar\'e duality
\cite[Theorem~3.30]{Hatcher2002}.  It applies to the
homology class pushed forward by $p$; the image of $p$ need not itself be
a submanifold.  Its image consists exactly of those $W$ with
$\dim\ker(W-I_d)\geq l$.

We recall the geometric interpretation of this class.  If $M$ is a
closed oriented manifold of dimension $l^2$ and $f:M\to\mathrm U(d)$,
then $\langle f^*\alpha(d,l),[M]\rangle$ is the intersection number of
$f$ and $p$.  When these maps are transverse, it is the signed count of
pairs $(x,(W,E))$ with $f(x)=p(W,E)$.  This follows by intersecting
$f\times p$ with the diagonal of $\mathrm U(d)\times\mathrm U(d)$:
transversality makes its inverse image a compact oriented zero-manifold,
whose signed count gives the intersection pairing.
In particular, disjoint images give intersection number zero: disjoint
maps from compact manifolds have disjoint sufficiently small transverse
perturbations.  The same interpretation holds for products of these
maps and classes.

The part of $\alpha(d,l)$ we need is
\begin{equation}\label{eq:fixed-incidence-leading-term}
 \alpha(d,l)=c_{d,l}e_{1,d}\cdots e_{l,d}
       +\text{terms with fewer than $l$ generators},
 \qquad c_{d,l}\ne0.
\end{equation}
For $l=0$, the map $p$ is the identity and we take
$\alpha(d,0)=1$ and $c_{d,0}=1$.
For $l>0$, a product of $u$ distinct generators has degree at least
$1+3+\cdots+(2u-1)=u^2$.  Since $\alpha(d,l)$ has degree $l^2$, its
expansion contains at most $l$ generators in each term, and the sole
possible term of length $l$ is the one displayed in
\eqref{eq:fixed-incidence-leading-term}.

To show that its coefficient is nonzero, consider
\[
 f:\mathrm U(l)\longrightarrow\mathrm U(d),
 \qquad f(X)=X\oplus(-I_{d-l}).
\]
There is exactly one intersection pair for $f$ and $p$: it has
$X=I_l$, $W=W_0=I_l\oplus(-I_{d-l})$, and $E=E_0=\mathbb C^l\oplus0$.
Indeed, the second summand has no fixed vectors, while an $l$-dimensional
fixed space in the first summand forces $X=I_l$.  This intersection is
transverse.  Identify the tangent space at $W_0$ with skew-Hermitian
matrices by right multiplication by $W_0^{-1}$.  Variations of $f$ fill
the upper-left block; variations of $p$ in the fiber over $E_0$ fill
the lower-right block.  For a linear map $Z:E_0\to E_0^\perp$, put
\[
 T=\begin{pmatrix}0&-Z^*\\ Z&0\end{pmatrix}.
\]
The path $(e^{tT}W_0e^{-tT},e^{tT}E_0)$ lies in $R(d,l)$ and its image
under $p$ has tangent vector
\[
 [T,W_0]W_0^{-1}=2T.
\]
These vectors supply every off-diagonal direction.  The unique
intersection therefore has multiplicity $1$ or $-1$.  The map $f$ is
homotopic to the coordinate inclusion, so
\eqref{eq:fixed-block-restriction} computes its pullback.  Terms with
an index above $l$ vanish; any remaining term with fewer than $l$
generators has degree less than $\dim\mathrm U(l)=l^2$.  Only the
displayed term in \eqref{eq:fixed-incidence-leading-term} can account for
the nonzero intersection number, proving $c_{d,l}\ne0$.
The same argument includes $l=d$, when the complementary blocks are
absent.

\smallskip\noindent
\emph{Computing the intersection number.}
We now apply these classes to the maps in the statement.  Write
$e_{a,k_i}^{(i)}$ for the generator pulled back from the $i$th factor of
$D$.  Primitivity, compatibility with block inclusions, and the inversion
formula give
\begin{equation}\label{eq:fixed-generator-pullback}
 F_j^*e_{a,d_j}
   =\sum_{i:\,k_i\geq a}b_{ij}e_{a,k_i}^{(i)}
       \qquad(1\leq a\leq d_j).
\end{equation}
To justify the use of multiplication here, write a block diagonal map
as the product of maps that act on one block and are the identity on
all others.  Each positive occurrence contributes one generator and
each negative occurrence its negative.  Right multiplication by $P_j$
is homotopic to the identity and has no effect on cohomology.

Consider the top-degree class
\[
 \Omega=\prod_{j=1}^nF_j^*\alpha(d_j,l_j)
       \in H^{\sum_i k_i^2}(D;\mathbb Q).
\]
The exterior algebra $H^*(D;\mathbb Q)$ has altogether $\sum_i k_i$
generators, all of positive degree.  A nonzero monomial of top degree
must contain every one of them.  By
\eqref{eq:fixed-generator-pullback}, pullback preserves the number of
generators in each term, unless that term vanishes.  It follows that
any use of a shorter term from
\eqref{eq:fixed-incidence-leading-term} contributes zero to $\Omega$.
Thus we need only compute the product of the displayed leading terms.

For each $a\in\{1,\ldots,k_1\}$, put
$t_a=|\{i:k_i\geq a\}|$.  The rows with $k_i\geq a$ are
$1,\ldots,t_a$, and the selected columns with $l_j\geq a$ are
$j_1,\ldots,j_{t_a}$.  Grouping the exterior factors of degree $2a-1$
together gives
\[
 \prod_{t=1}^{t_a}
   \left(\sum_{i=1}^{t_a}b_{i,j_t}e_{a,k_i}^{(i)}\right)
 =\det B[1,\ldots,t_a\mid j_1,\ldots,j_{t_a}]
       \prod_{i=1}^{t_a}e_{a,k_i}^{(i)}.
\]
Consequently the coefficient of the top monomial in $\Omega$, up to
an overall sign, is
\begin{equation}\label{eq:fixed-top-coefficient}
 \left(\prod_{j=1}^n c_{d_j,l_j}\right)
 \prod_{a=1}^{k_1}
   \det B[1,\ldots,t_a\mid j_1,\ldots,j_{t_a}].
\end{equation}
Every factor is nonzero, by
\eqref{eq:fixed-nested-minors} and
\eqref{eq:fixed-incidence-leading-term}.  Hence
$\langle\Omega,[D]\rangle\ne0$.

Finally, let $F=(F_1,\ldots,F_n)$ and let
$p_{\mathrm{tot}}=\prod_j p$ be the map from
$\prod_jR(d_j,l_j)$ to $\prod_j\mathrm U(d_j)$.
Up to the harmless orientation sign, $\Omega$ is the pullback by $F$
of the Poincar\'e dual of the cycle defined by $p_{\mathrm{tot}}$.
Its nonzero evaluation implies that these two maps have intersecting
images.  At an intersection, $F_j(X)$ fixes an $l_j$-plane for every
$j$, which gives \eqref{eq:fixed-space-bound} by
\eqref{eq:fixed-target-dimensions}.
\end{proof}
